\section*{Erd\H{o}s Problem 992: a metric discrepancy counterexample}
\subsection*{Statement and attribution}
The proposed bounds $\sqrt N(\log N)^{o(1)}$ and
$\sqrt N(\log\log N)^{O(1)}$ are false for general increasing integer sequences.
We give a complete block construction proving, for one fixed sequence and almost every $\alpha$,
\[
 \limsup_{N\to\infty}\frac{D(N;\alpha)}{\sqrt{N\log N}}=\infty.
\]
The known negative result is due to Berkes and Philipp, \emph{The size of trigonometric and Walsh series and uniform distribution mod 1}, J. London Math. Soc. 50 (1994), 454--464. The proof below reconstructs the independent-digit block mechanism; no novelty is claimed.

\subsection*{Blocks and independent exceptional events}
For integers $i\ge3$, put
\[
 r_i=\lceil i\log i\log\log i\rceil,\qquad
 \ell_i=\lceil\log_2(6r_i)\rceil.
\]
Choose $b_3=0$ and $b_{i+1}=b_i+\ell_i+1$, and set $M_i=2^{b_i}$.
Our increasing sequence is the concatenation of the blocks
\[
 M_i,\,2M_i,\ldots,r_iM_i\qquad(i=3,4,\ldots).
\]
There is no overlap or reordering issue: $M_{i+1}=2^{\ell_i+1}M_i>r_iM_i$.

Except for the null set of dyadic rationals, the binary digits of a Lebesgue-uniform $\alpha\in[0,1)$ are independent fair bits. Let $E_i$ be the event that digits
$b_i+1,\ldots,b_i+\ell_i$ are all zero.
These digit ranges are disjoint, so the events are independent. Moreover
\[
 \frac1{12r_i}<\lambda(E_i)=2^{-\ell_i}\le\frac1{6r_i},
 \qquad
 \sum_{i\ge3}\frac1{r_i}=\infty.
\]
The divergence follows from the integral of $1/(x\log x\log\log x)$; rounding only changes a constant factor.
The independent Borel--Cantelli lemma therefore gives infinitely many $E_i$ for almost every $\alpha$.

On $E_i$, $0\le\{M_i\alpha\}<2^{-\ell_i}\le1/(6r_i)$.
Thus for every $1\le j\le r_i$,
$\cos(2\pi jM_i\alpha)\ge1/2$, and the whole block sum is at least $r_i/2$.

\subsection*{A large block forces a large prefix}
Let $N_i=\sum_{j=3}^{i}r_j$ and
$S_N(\alpha)=\sum_{n\le N}\cos(2\pi\alpha x_n)$ for the constructed sequence $(x_n)$.
On $E_i$,
\[
 S_{N_i}-S_{N_{i-1}}\ge r_i/2,
\]
so at least one of $|S_{N_i}|,|S_{N_{i-1}}|$ is at least $r_i/4$.
Both possible indices tend to infinity with $i$. For all large $i$,
\[
 N_i\ll i^2\log i\log\log i,\qquad \log N_i\ll\log i.
\]
Consequently, at one of these two endpoints $N$,
\[
 \frac{|S_N|}{\sqrt{N\log N}}
 \ge\frac{r_i/4}{\sqrt{N_i\log N_i}}
 \gg\sqrt{\log\log i}\longrightarrow\infty
\]
along the infinitely many successful blocks.
This step permits earlier blocks to have either sign; it does not assume that their sums are negligible.

\subsection*{From a cosine sum to interval discrepancy}
For points $y_1,\ldots,y_N\in[0,1)$ define
$F(t)=\#\{j:y_j\le t\}-Nt$.
Integration by parts against the signed measure $\sum_j\delta_{y_j}-N\,dt$, of total mass zero, gives
\[
 \left|\sum_{j=1}^{N}\cos(2\pi y_j)\right|
 \le\sup_t|F(t)|\int_0^1|(\cos(2\pi t))'|\,dt
 =4\sup_t|F(t)|\le4D_N.
\]
Endpoint conventions do not affect the supremum; alternatively use limits of half-open intervals.
Apply this with $y_j=\{\alpha x_j\}$. The preceding prefix estimate proves the displayed discrepancy limsup.

\subsection*{Conclusion and comparison with the conjectures}
For any fixed $C$, $\sqrt{N}(\log\log N)^C=o(\sqrt{N\log N})$.
Also $\sqrt{N}(\log N)^{o(1)}=o(\sqrt{N\log N})$. Our single integer sequence therefore disproves both assertions on a set of $\alpha$ of full measure.
The proof uses only independence of disjoint binary digits, Borel--Cantelli, and the displayed elementary discrepancy inequality. It does not use an invalid uniform truncation estimate for the discontinuous sawtooth function.
Sources: \url{https://www.erdosproblems.com/992} and \url{https://doi.org/10.1112/jlms/50.3.454}.
\subsection*{COMPLETION ESTIMATE}
\textbf{COMPLETION: 100\%}
